\appendix
\chapter[Appendix]{Appendix}

\section{Supplementary results}

\subsection{Asymptotics of recurrences that include a sum of history}
\label{subsec:asymptotics-of-recurrences-that-include-a-sum-of-history}

\begin{proposition}
  \label{lem:full-history-sum-asymptotics}
  Consider the sequence $(z_n)$ defined by the recurrence
  \begin{equation}
    \label{eq:full-history-sum-recurrence}
    z_{n+2} = \dfrac{K}{n}\sum\limits_{i=1}^n z_{i}, \quad n \ge 1,
  \end{equation}
  where $K > 0$. The initial values $z_1$ and $z_2$ are nonnegative and not both zero. Then $z_n\sim Cn^{K-1}$ for some constant $C > 0$.
\end{proposition}

\begin{proof}
  The recurrence is equivalent to
  $$Kz_{n+2} = \sum_{i=1}^n z_i.$$
  Let $Z(x) = \sum_{n=1}^\infty z_n x^n$ be the generating function of $(z_n)$. Summing the rewritten recurrence over $n\ge 1$ gives
  $$\sum_{n=1}^\infty n z_{n+2} x^n = K \sum_{n=1}^\infty \sum_{n=1}^\infty z_n x^n.$$

  For the left-hand side, we have
  \begin{align*}
    \sum_{n=1}^\infty n z_{n+2} x^n
     & = \frac{1}{x} \sum_{n=3}^\infty (n-2) z_n x^{n-1} = \frac{1}{x} \left( \sum_{k=3}^\infty k z_k x^{k-1} - \frac{2}{x} \sum_{k=3}^\infty z_k x^k \right) \\
     & = \frac{1}{x} \left(Z'(x) - z_1 - 2z_2 x - \frac{2}{x} \left(Z(x) - z_1 x - z_2 x^2\right) \right)                                                     \\
     & = \frac{Z'(x)}{x} - \frac{2Z(x)}{x^2} + \frac{z_1}{x}.
  \end{align*}
  For the right-hand side, we have
  $$K \sum_{n=1}^\infty \sum_{n=1}^\infty z_n x^n = K \frac{Z(x)}{1-x}.$$
  Equating the two sides and rearranging, we obtain a first-order linear differential equation
  $$Z'(x) - \left( \frac{2}{x} + \frac{Kx}{1-x} \right) Z(x) = -z_1.$$
  Multiplying both sides by the integrating factor $I(x) = x^{-2} (1-x)^K e^{Kx}$ gives
  $$[I(x) Z(x)]' = -z_1 x^{-2} (1-x)^K e^{Kx} = -z_1 x^{-2} - z_1 \frac{(1-x)^K e^{Kx} - 1}{x^2}.$$
  To integrate this from $0$ to $x$ without introducing an arbitrary integration constant, we group the singular $x^{-2}$ terms. The local expansion of $Z(x)$ at $x=0$ is $z_1 x + z_2 x^2 + \mathcal{O}(x^3)$, and the Taylor series of $(1-x)^K e^{Kx}$ begins with $1 - \frac{K}{2}x^2 + \mathcal{O}(x^3)$. Thus, we have the limit
  $$\lim_{x\to 0} \left( I(x)Z(x) - \frac{z_1}{x} \right) = \lim_{x\to 0} \left( x^{-2} (1-x)^K e^{Kx} Z(x) - \frac{z_1}{x} \right) = z_2.$$
  Furthermore, $\lim_{x\to 0} \frac{(1-x)^K e^{Kx} - 1}{x^2} = -\frac{K}{2}$, confirming that the right-hand integrand is continuous. Moving $-z_1 x^{-2}$ to the left side and integrating from $0$ to $x$ via the Fundamental Theorem of Calculus yields
  $$I(x) Z(x) - \frac{z_1}{x} - z_2 = -z_1 \int_0^x \frac{(1-t)^K e^{Kt} - 1}{t^2} \, \mathrm{d}t.$$
  Multiplying by $I(x)^{-1}$ isolates the closed-form generating function
  $$Z(x) = (1-x)^{-K} e^{-Kx} \left( z_2 x^2 + z_1 x - z_1 x^2 \int_0^x \frac{(1-t)^K e^{Kt} - 1}{t^2} \, \mathrm{d}t \right).$$
  To extract the asymptotic behavior of $(z_n)$, we analyze $Z(x)$ near its dominant singularity at $x=1$. The integral converges on $(0, 1)$, giving the local expansion
  $$Z(x) \underset{x \to 1}{\sim} A e^{-K} (1-x)^{-K},$$
  where the leading constant is
  $$A = z_2 + z_1 \left( 1 - \int_0^1 \frac{(1-t)^K e^{Kt} - 1}{t^2} \, \mathrm{d}t \right).$$
  We must ensure singularity cancellation does not occur. For any $t \in (0, 1)$, the inequality $\ln(1-t) < -t$ implies $K\ln(1-t) + Kt < 0$. Exponentiating yields $(1-t)^K e^{Kt} < 1$, so the integrand is strictly negative. Consequently, the integral evaluates to a strictly negative value, making $1 - \int_0^1 \frac{(1-t)^K e^{Kt} - 1}{t^2} \, \mathrm{d}t > 1$. Because $z_1, z_2 \ge 0$ and are not both zero, it follows that $A \ge z_2 + z_1 > 0$. Since $K > 0$, we apply Theorem \ref{thm:sim-transfer}. Mapping the $(1-x)^{-K}$ singularity to the coefficients $z_n = [x^n]Z(x)$ yields
  $$z_n \sim \frac{A e^{-K}}{\Gamma(K)} n^{K-1}.$$
  Defining $C = \frac{A e^{-K}}{\Gamma(K)}$, we have $C > 0$, completing the proof.
\end{proof}

\subsection{Spaces of probability measures}
\label{subsec:supp-locally-convex-spaces}

\begin{definition}[Weak convergence of measures]
  \label{def:weak-convergence}
  Let $(M,d)$ be a metric space and let $\M$ be the Borel $\sigma$-algebra on $M$. Let a probability distribution be defined on $(M, \M)$, and let $C_b(M)$ be the set of real-valued, bounded, continuous functions on $M$. A sequence of measures $(\mu_n)$ converges weakly to a measures $\mu$ if, for all $f\in C_b(M)$, we have
  $$\int f \, d\mu_n \to \int f \, d\mu.$$
\end{definition}

\begin{definition}[Tightness]
  \label{def:tight}
  A sequence of probability measures $(\mu_n)$ on $\RR$ is tight if for every $\varepsilon > 0$, there exists $M > 0$ such that $\sup_n \mu_n(\{x: |x| > M\}) < \varepsilon$.

  A sequence of random variables $(Z_n)$ is tight if for every $\varepsilon > 0$, there exists $M > 0$ such that $\sup_n\PP(|Z_n| > M) < \varepsilon$.
\end{definition}

\begin{theorem}[Prokhorov's theorem]
  \label{thm:prokhorov}
  A sequence of probability measures $(\mu_n)$ on $\RR$ is tight if and only if it is relatively compact in the weak topology (i.e., its closure is compact).
\end{theorem}

\section{Remaining proofs}

\subsection{Lemma \ref{lem:limiting-equation}}
\label{subsec:proof-limiting-equation}

We need the following lemma.

\begin{lemma}
  \label{lemma:average_uniform_convergence}
  Let $U_n$ be a uniform random variable on $\{1, 2, \ldots, n\}$ and let $U$ be a uniform random variable on $[0,1]$. Then $\dfrac{U_n}{n}\xlongrightarrow{d} U$.
\end{lemma}

\begin{proof}
  Recall the CDF of $U$
  $$F_U(x) = \begin{cases}
      0 & x < 0,         \\
      x & 0 \le x \le 1, \\
      1 & x > 1.
    \end{cases}$$
  Let $V_n = U_n/n$. We have $F_{V_n}(x) = \PP(U_n/n \le x) = \PP(U_n\le nx)$. Hence
  \begin{itemize}
    \item If $x\le 0$, then $F_{V_n}(x) = 0 = F_U(x)$, since $U_n\ge 1$.
    \item If $x > 1$, then $F_{V_n}(x) = 1 = F_U(x)$, since $U_n\le n$.
    \item If $0 < x \le 1$, then $F_{V_n}(x) = \dfrac{\lfloor nx \rfloor}{n}$. Since $x - \dfrac{1}{n} \le \dfrac{\lfloor nx \rfloor}{n} \le x$, we have $F_{V_n}(x) \to x = F_U(x)$.
  \end{itemize}
\end{proof}


\begin{proof}[Proof of Lemma \ref{lem:limiting-equation}]
  We first show that
  $$\left(\dfrac{U_n}{n+2}\right)^\alpha X_{U_n}^{(1)} + \left(\dfrac{n+1-U_n}{n+2}\right)^\alpha X_{n+1-U_n}^{(2)} \xlongrightarrow{d} U^\alpha X^{(1)} + (1-U)^\alpha X^{(2)}.$$

  By Lemma \ref{lemma:average_uniform_convergence} and Slutsky's theorem,
  $$\dfrac{U_n}{n+2} = \dfrac{U_n}{n} \cdot \dfrac{n}{n+2} \xlongrightarrow{d} U.$$

  To apply the Continuous Mapping Theorem, we need to show that
  $$\left(\dfrac{U_n}{n+2}, X_{U_n}^{(1)}, X_{n+1-U_n}^{(2)}\right) \xlongrightarrow{d} \left(U, X^{(1)}, X^{(2)}\right).$$

  Since $U, X^{(1)}, X^{(2)}$ are independent, this is equivalent to showing that $\forall u, x_1, x_2 \in \mathbb{R}$
  $$\PP\left(U_n \le \lfloor nu \rfloor, X_{U_n}^{(1)} \le x_1, X_{n+1-U_n}^{(2)} \le x_2\right) \longrightarrow \PP(U \le u) \PP(X^{(1)} \le x_1) \PP(X^{(2)} \le x_2).$$

  Indeed,
  \begin{align*}
     & \PP\left(U_n \le \lfloor nu \rfloor, X_{U_n}^{(1)} \le x_1, X_{n+1-U_n}^{(2)} \le x_2\right)                                                                         \\
     & = \sum_{k=1}^{\lfloor nu \rfloor} \PP(U_n=k) \PP\left(X_{k}^{(1)} \le x_1, X_{n+1-k}^{(2)} \le x_2 | U_n=k\right)                                                    \\
     & =  \dfrac{1}{n}\sum_{k=1}^{\lfloor nu \rfloor} \PP\left(X_{k}^{(1)} \le x_1, X_{n+1-k}^{(2)} \le x_2\right)               & (U_n\bot (X_{k}^{(1)}, X_{n+1-k}^{(2)})) \\
     & =  \dfrac{1}{n}\sum_{k=1}^{\lfloor nu \rfloor} \PP\left(X_{k}^{(1)} \le x_1\right)\PP\left(X_{n+1-k}^{(2)} \le x_2\right) & (X_{k}^{(1)} \bot X_{n+1-k}^{(2)})       \\
  \end{align*}

  Let $F_n$ be the common CDF of $X_n^{(1)}$ and $X_n^{(2)}$, and $F$ be the CDF of $X^{(1)}$. Let $m = \lfloor nu \rfloor$. The problem becomes showing that
  $$\dfrac{1}{n}\sum_{k=1}^{m} F_k(x_1)F_{n+1-k}(x_2) \longrightarrow u F(x_1) F(x_2).$$
  Equivalently,
  $$\dfrac{m}{n}\left(\dfrac{1}{m}\sum_{k=1}^{m} \left(F_k(x_1)F_{n+1-k}(x_2)-F(x_1) F(x_2)\right)\right) \longrightarrow 0.$$

  We have
  \begin{align*}
     & \dfrac{1}{m}\sum_{k=1}^{m} \left|F_k(x_1)F_{n+1-k}(x_2)-F(x_1) F(x_2)\right|                                                                                                       \\
     & \le \dfrac{1}{m}\sum_{k=1}^{m} \left|F_k(x_1)\left(F_{n+1-k}(x_2)-F(x_2)\right)\right| + \dfrac{1}{m}\sum_{k=1}^{m} \left|F(x_2)\left(F_k(x_1)-F(x_1)\right)\right|                \\
     & \le \dfrac{1}{m}\sum_{k=1}^{m} \left|F_{n+1-k}(x_2)-F(x_2)\right| + \dfrac{1}{m}\sum_{k=1}^{m} \left|F_k(x_1)-F(x_1)\right|                                         & (|F_n|\le 1) \\
  \end{align*}
  The second term goes to zero by the fact that $F_n(x_1)\to F(x_1)$ and Cesàro Means. We can write the first term as
  \begin{align*}
    \dfrac{1}{m}\sum_{k=n-m+1}^{n} \left|F_{k}(x_2)-F(x_2)\right|.
  \end{align*}

  If $u=1$, its convergence is similar to the second term. If $u<1$, then $n-m+1 \to \infty$. By pointwise convergence, $\forall \epsilon>0, \exists N, \forall n\ge N: |F_n(x_2)-F(x_2)|<\epsilon$. Hence, for $n-m+1\ge N$,
  \begin{align*}
    \dfrac{1}{m}\sum_{k=n-m+1}^{n} \left|F_{k}(x_2)-F(x_2)\right| < \epsilon.
  \end{align*}

  Using Slutsky's theorem again, we have

  $$\left(\dfrac{n+1}{n+2}, \dfrac{U_n}{n+2}, X_{U_n}^{(1)}, X_{n+1-U_n}^{(2)}\right) \xlongrightarrow{d} \left(1, U, X^{(1)}, X^{(2)}\right).$$

  Now define the function $f:\mathbb{R}^4_+\cap \left\{(t,u,x,y): t\ge u\right\} \to \mathbb{R}$ by
  $$f(t,u,x,y) = u^\alpha x + (t-u)^\alpha y.$$
  Then $f$ is continuous. The Continuous Mapping Theorem applies. The convergence
  $$\max\left(\left(\dfrac{U_n}{n+2}\right)^\alpha X_{U_n}^{(1)}, \left(\dfrac{n+1-U_n}{n+2}\right)^\alpha X_{n+1-U_n}^{(2)}\right) \xlongrightarrow{d} \max\left(U^\alpha X^{(1)}, (1-U)^\alpha X^{(2)}\right)$$

  is similar. And finally, we define the continuous function $g(i,j,u,v) = iu+jv$ on $\RR^4$ to finish the proof.
\end{proof}


\subsection{Proposition \ref{prop:f}}
\label{subsec:proof-prop-f}

\begin{proof}
  The function $f$ is strictly decreasing in $\alpha$. Next, we show that $f$ is continuous. The first term is continuous. We show that the second term is Lipschitz continuous. Let $\alpha, \alpha' \in (0,1]$ and assume without loss of generality that $\alpha > \alpha'$. We have
  \begin{align*}
     & \EE\left[\max\left(U^\alpha X^{(1)}, (1-U)^\alpha X^{(2)}\right) - \max\left(U^{\alpha'} X^{(1)}, (1-U)^{\alpha'} X^{(2)}\right)\right] \\
     & \le \EE\left[\max\left(\left|U^\alpha - U^{\alpha'}\right|X^{(1)}, \left|(1-U)^\alpha - (1-U)^{\alpha'}\right|X^{(2)}\right)\right]     \\
     & \le \EE\left[\left|U^\alpha - U^{\alpha'}\right|\right] + \EE\left[\left|(1-U)^\alpha - (1-U)^{\alpha'}\right|\right]                   \\
     & = 2\EE\left[\left|U^\alpha - U^{\alpha'}\right|\right]
     & \text{($U \stackrel{d}{=} 1-U$)}                                                                                                        \\
     & = 2\EE\left[U^{\alpha'} - U^{\alpha}\right]
     & \text{($U^\alpha\ge U^{\alpha'}, a.s.$)}                                                                                                \\
     & = \dfrac{2}{\alpha'+1} - \dfrac{2}{\alpha+1} = \dfrac{2(\alpha-\alpha')}{(\alpha+1)(\alpha'+1)}                                         \\
     & \le 2(\alpha-\alpha').
  \end{align*}

  Finally, we evaluate $f$ at the endpoints.
  $$\lim_{\alpha \to 0+} f(\mu, \alpha) = 2u+v\EE\left[\max\left(X^{(1)}, X^{(2)}\right)\right] > 2u+v > 1,$$
  $$f(1) = u + v\EE\left[\max\left(U^\alpha X^{(1)}, (1-U)^\alpha X^{(2)}\right)\right] \le u+v \le 1.$$
\end{proof}

\subsection{Proposition \ref{prop:uniform-bounded-second-moment}}
\label{subsec:proof-prop-phi}

\begin{proof}
  Let $\mu\in\D_2(1)$ have finite second moment, and let $Z\sim \mu$. Following the steps in the previous section, we have
  \begin{align*}
    \EE\left[\Phi(Z)^2\right]
     & \le \dfrac{2(u+v)}{2\alpha+1}\EE\left[Z^2\right] + 2uB(\alpha+1, \alpha+1) \\
     & < \dfrac{2(u+v)}{2(2u+v)-1}\EE\left[Z^2\right] + 2uB(2u+v, 2u+v)           \\
  \end{align*}
  Choose
  $$M_0 = 2uB(2u+v, 2u+v)\dfrac{2(2u+v)-1}{2u-1}.$$
  We have $\EE\left[\Phi(Z)^2\right] \le M$. Hence, $\Phi$ maps $\D_2(1, M)$ into itself.
\end{proof}

\subsection{Proposition \ref{prop:nonempty-compact-convex}}
\label{subsec:proof-prop-locally-convex-space}

\begin{proof}
  The set $\D_2(1, M)$ is nonempty since it contains $\delta_1$.

  For convexity, let $\mu, \nu \in \D_2(1, M)$. Then for every $\lambda \in [0,1]$, we have
  \begin{align*}
    \int_0^\infty x d(\lambda \mu + (1-\lambda)\nu)
     & = \lambda \int_0^\infty x d\mu + (1-\lambda) \int_0^\infty x d\nu = 1,       \\
    \int_0^\infty x^2 d(\lambda \mu + (1-\lambda)\nu)
     & = \lambda \int_0^\infty x^2 d\mu + (1-\lambda) \int_0^\infty x^2 d\nu \le M.
  \end{align*}

  For compactness, we use Prokhorov's theorem to reduce to proving that $\D_2(1, M)$ is tight and closed. Tightness is guaranteed by Chebyshev's inequality: for any $\mu \in \D_2(1, M)$ and any $R > 0$

  \begin{equation}
    \mu(\{x \ge R\}) \le \dfrac{1}{R^2}\int_{0}^\infty x^2 \,\mathrm{d}\mu \le \dfrac{M}{R^2}.
  \end{equation}

  For closedness, let $(\mu_n)$ be a sequence in $\D_2(1, M)$ that converges weakly to $\mu$. We show that $\mu \in \D_2(1, M)$. The boundedness of the second moment follows from a generalized Fatou's lemma \cite[Theorem 1.1]{feinberg2014fatou}
  $$\int_0^\infty x^2 \,\mathrm{d}\mu \le \liminf_{n \to \infty} \int_0^\infty x^2 \,\mathrm{d}\mu_n \le M.$$

  The function $x\mapsto x$ is uniformly integrable with respect to $\D_2(1, M)$ as follows. For any $\mu \in \D_2(1, M)$ and any $R > 0$
  $$\int_{R}^\infty x \,\mathrm{d}\mu \le \int_{R}^\infty x \left(\frac{x}{R} \right) \mathrm{d}\mu = \frac{1}{R} \int_{R}^\infty x^2 \,\mathrm{d}\mu \le \frac{1}{R} \int_{0}^\infty x^2 \,\mathrm{d}\mu \le \frac{M}{R}.$$
  which tends to 0 as $R \to \infty$. Therefore, the first moment converges \cite[Theorem 3.5]{billingsley1999convergence}
  $$\int_0^\infty x \,\mathrm{d}\mu_n \to \int_0^\infty x \,\mathrm{d}\mu.$$
\end{proof}

\subsection{Proposition \ref{prop:phi-continuity}}
\label{subsec:proof-prop-phi-continuity}

\begin{lemma}
  \label{lem:collapse-one-step}
  For every $\mu,\nu\in \D_2(1)$ and $\alpha\in(0,1]$, we have
  $$d_2(T_\alpha(\mu), T_\alpha(\nu)) \le \sqrt{\dfrac{2(u+v)}{2\alpha+1}}d_2(\mu, \nu) \text{ and } d_1(T_\alpha(\mu), T_\alpha(\nu)) \le \dfrac{2(u+v)}{\alpha+1}d_1(\mu, \nu).$$
\end{lemma}

\begin{proof}
  Since $I$ and $J$ are Bernoulli, we have $I^2 = I$ and $J^2 = J$. Also, from $I+J\le 1$, we have $IJ = 0$. Therefore, for every $X,X^{(1)},X^{(2)}\stackrel{\text{iid}}{\sim}\mu$ and $Y,Y^{(1)},Y^{(2)}\stackrel{\text{iid}}{\sim}\nu$, we have
  \begin{align*}
    d_2^2(T_\alpha(\mu), T_\alpha(\nu))
     & \le \EE\left[I\left(\left(U^\alpha\left(X^{(1)}-Y^{(1)}\right) + \left(1-U\right)^\alpha\left(X^{(2)}-Y^{(2)}\right)\right)\right)^2\right]                                                      \\
     & \,\,\,\,\,+\EE \left[J\left.\left(\max\left(U^\alpha X^{(1)}, \left(1-U\right)^\alpha X^{(2)}\right) - \max\left(U^\alpha Y^{(1)}, \left(1-U\right)^\alpha Y^{(2)}\right)\right)\right)^2\right] \\
     & := A+B.
  \end{align*}
  For the first expectation $A$, we use $\EE[X] = \EE[Y] = 1$ and get
  \begin{align*}
    A
     & = u\EE\left[U^{2\alpha}\left(X^{(1)}-Y^{(1)}\right)^2 + (1-U)^{2\alpha}\left(X^{(2)}-Y^{(2)}\right)^2 \right] \\
     & = \dfrac{2u}{2\alpha+1}\EE[(X-Y)^2].
  \end{align*}
  For the second one, $B$, we use the inequality
  $$(\max(a,b) - \max(c,d))^2 \le \max(a-c, b-d)^2 \le (a-c)^2+(b-d)^2$$
  for real numbers $a,b,c,d$ and get
  \begin{align*}
    B
     & \le v\EE\left[U^{2\alpha}\left(X^{(1)}-Y^{(1)}\right)^2 + (1-U)^{2\alpha}\left(X^{(2)}-Y^{(2)}\right)^2 \right] \\
     & = \dfrac{2v}{2\alpha+1}\EE[(X-Y)^2].
  \end{align*}
  Therefore,
  $$d_2^2(T_\alpha(\mu), T_\alpha(\nu)) \le \dfrac{2(u+v)}{2\alpha+1}\EE[(X-Y)^2].$$
  Taking the expectation over $X\sim\mu$ and $Y\sim\nu$, we get the desired result. The case for $d_1$ is similar.
\end{proof}

\begin{proof}[Proof of Proposition \ref{prop:phi-continuity}]
  Let $(\mu_n)$ be a sequence in $\D_2(1, M)$ that converges weakly to $\mu$, which is equivalent to $d_1(\mu_n, \mu) \to 0$. To show that $\Phi(\mu_n)\stackrel{w}{\longrightarrow}\Phi(\mu)$, we will show that $d_1(\Phi(\mu_n), \Phi(\mu)) \to 0$ by showing that $\Phi$ is Lipschitz with respect to $d_1$. Let $\mu, \nu \in \D_2(1, M)$. Write $\alpha_1 := \alpha(\mu)$ and $\alpha_2 := \alpha(\nu)$. Without loss of generality, assume that $\alpha_1 > \alpha_2$. We have
  \begin{align*}
    d_1(\Phi(\mu), \Phi(\nu))
     & = d_1(T_{\alpha_1}(\mu), T_{\alpha_2}(\nu))                                                \\
     & \le d_1(T_{\alpha_2}(\mu), T_{\alpha_2}(\nu)) + d_1(T_{\alpha_1}(\mu), T_{\alpha_2}(\mu)).
  \end{align*}
  By Lemma~\ref{lem:collapse-one-step}, we have
  $$d_1(T_{\alpha_2}(\mu), T_{\alpha_2}(\nu)) \le \dfrac{2(u+v)}{\alpha_2+1}d_1(\mu, \nu) \le 2(u+v) d_1(\mu, \nu).$$
  For the second distance, we let $Y^{(1)},Y^{(2)} \sim \nu$. Then, following the same steps as in Lemma~\ref{prop:f}, we have
  \begin{align*}
    d_1(T_{\alpha_1}(\mu), T_{\alpha_2}(\mu))
     & \le 2(u+v)(\alpha_1 - \alpha_2).
  \end{align*}
  We have $f(\mu, \alpha_1) = f(\nu, \alpha_2) = 1$. Hence
  $$|f(\mu, \alpha_1) - f(\mu, \alpha_2)| = |f(\mu, \alpha_2) - f(\nu, \alpha_2)|.$$
  Consider the left-hand side term where we fix the distribution $\mu$. We rewrite $f$ as
  $$f(\mu, \alpha) = \dfrac{2u}{\alpha+1} + v\EE\left[U^\alpha Z^{(1)}\bm{1}_{A} + (1-U)^\alpha Z^{(2)}\bm{1}_{A^c}\right], \quad Z^{(1)}, Z^{(2)} \stackrel{\text{iid}}{\sim} \mu,$$
  where $A$ is the event that $U^\alpha Z^{(1)} \ge (1-U)^\alpha Z^{(2)}$. Since $\PP\left(U^\alpha Z^{(1)} = (1-U)^\alpha Z^{(2)}\right) = 0$, we have
  \begin{align*}
    \dfrac{\partial}{\partial \alpha} f(\mu, \alpha)
     & = -\dfrac{2u}{(\alpha+1)^2} + v\EE\left[\ln(U) U^\alpha Z^{(1)} \mathbf{1}_A + \ln(1-U) (1-U)^\alpha Z^{(2)} \mathbf{1}_{A^c}\right].
  \end{align*}
  Since $\ln(U) < 0$ and $\ln(1-U) < 0$, we have for every $\alpha\in(0,1]$,
  $$\left|\dfrac{\partial}{\partial \alpha} f(\mu, \alpha)\right| \ge \dfrac{2u}{(\alpha+1)^2}.$$
  By Mean Value Theorem, we have
  $$|f(\mu, \alpha_1) - f(\mu, \alpha_2)| \ge \dfrac{2u}{(\alpha_1+1)^2}|\alpha_1 - \alpha_2|\ge \dfrac{u}{2}|\alpha_1 - \alpha_2|$$
  We also have
  $$|f(\mu, \alpha_2) - f(\nu, \alpha_2)| \le \dfrac{2v}{\alpha_2+1}d_1(\mu, \nu) \le 2v d_1(\mu, \nu).$$
  Thus,
  $$d_1(\Phi(\mu), \Phi(\nu)) \le 2(u+v)\left(1+\dfrac{4v}{u}\right)d_1(\mu, \nu).$$
\end{proof}
